\section{System Architecture}
\label{sec:architecture}

\subsection{Design Principles}

Miniforge is designed around five core principles that guide its architecture and implementation:

\begin{enumerate}
    \item \textbf{Async-First:} All operations are asynchronous by default, enabling concurrent processing and non-blocking APIs.
    \item \textbf{Hardware-Aware:} System behavior adapts to available memory, CPU cores, and platform characteristics.
    \item \textbf{Backend-Agnostic:} Clean abstraction over multiple inference backends with automatic fallback.
    \item \textbf{Memory-Conservative:} Aggressive memory management with safety margins and automatic optimization.
    \item \textbf{Quality-Preserving:} Quantization and compression decisions prioritize maintaining model capabilities.
\end{enumerate}

\subsection{High-Level Architecture}

Figure \ref{fig:architecture} presents the high-level system architecture.

\begin{figure}[h]
\centering
\begin{tikzpicture}[
    node distance=1.5cm,
    box/.style={rectangle, draw, minimum width=3cm, minimum height=0.8cm, align=center},
    arrow/.style={->, >=stealth, thick}
]
    % Input layer
    \node[box, fill=blue!20] (api) {API Layer\\(Async Interface)};
    
    % Core layer
    \node[box, fill=green!20, below of=api] (engine) {Inference Engine\\(Unified Backend)};
    
    % Backend layer
    \node[box, fill=yellow!20, below left=1.5cm and 0.5cm of engine] (llama) {llama.cpp\\Backend};
    \node[box, fill=yellow!20, below right=1.5cm and 0.5cm of engine] (transformers) {Transformers\\Backend};
    
    % Memory management
    \node[box, fill=red!20, right=3cm of engine] (memory) {Memory\\Manager};
    
    % Configuration
    \node[box, fill=purple!20, left=3cm of engine] (config) {Configuration\\(M7Config)};
    
    % Arrows
    \draw[arrow] (api) -- (engine);
    \draw[arrow] (engine) -- (llama);
    \draw[arrow] (engine) -- (transformers);
    \draw[arrow, <->] (engine) -- (memory);
    \draw[arrow, <->] (engine) -- (config);
    \draw[arrow, <->] (memory) -- (config);
    
    % Additional modules
    \node[box, fill=orange!20, below=3cm of engine] (tools) {Tool\\Executor};
    \node[box, fill=orange!20, right=0.5cm of tools] (vision) {Vision\\Processor};
    \node[box, fill=orange!20, left=0.5cm of tools] (streaming) {Streaming\\Handler};
    
    \draw[arrow, <->] (engine) -- (tools);
    \draw[arrow, <->] (engine) -- (vision);
    \draw[arrow, <->] (engine) -- (streaming);
    \draw[arrow, <->] (memory) -- (llama);
    \draw[arrow, <->] (memory) -- (transformers);
\end{tikzpicture}
\caption{Miniforge System Architecture}
\label{fig:architecture}
\end{figure}

The architecture consists of five layers:

\subsubsection{API Layer}

Provides user-facing interfaces for chat, completion, streaming, tool calling, and vision processing. All APIs are async-first with synchronous wrappers for convenience.

\subsubsection{Inference Engine}

The core abstraction layer managing backend initialization, request routing, and result aggregation. Implements the state machine for model lifecycle (loading $\rightarrow$ ready $\rightarrow$ generating $\rightarrow$ cleanup).

\subsubsection{Backend Layer}

Pluggable inference backends with shared interface. Currently implements:

\begin{itemize}
    \item \textbf{llama.cpp:} Primary backend for CPU-optimized GGUF inference
    \item \textbf{Transformers:} Fallback for native HuggingFace model support
\end{itemize}

\subsubsection{Memory Management}

Monitors system memory, tracks model and cache allocations, enforces limits, and triggers optimizations when approaching constraints.

\subsubsection{Configuration}

Hardware-specific configuration management with profiles for different platforms. The M7Config class encapsulates optimal settings for the GMKtech M7.

\subsection{Module Interactions}

The interaction flow for a typical chat request follows this sequence:

\begin{enumerate}
    \item User calls \texttt{model.chat()} with message and optional parameters
    \item API layer validates and normalizes input
    \item InferenceEngine formats messages into prompt template
    \item MemoryManager checks available resources
    \item Backend initializes if not already loaded
    \item Generation proceeds with optional streaming
    \item Results returned through API layer
\end{enumerate}

For streaming requests, tokens are yielded through an async iterator as they become available, rather than accumulating the complete response.

\subsection{Component Specifications}

\subsubsection{Miniforge Class}

The main user-facing interface encapsulates model management:

\begin{lstlisting}[language=Python, caption=Miniforge Interface]
class Miniforge:
    async def chat(
        self,
        message: str,
        history: Optional[List[Dict]] = None,
        system_prompt: Optional[str] = None,
        tools: Optional[List[Tool]] = None,
        stream: bool = False,
        **generation_params
    ) -> Union[str, AsyncIterator[str]]
    
    async def generate(
        self,
        prompt: str,
        max_tokens: int = 512,
        temperature: float = 1.0,
        stream: bool = False,
    ) -> Union[str, AsyncIterator[str]]
    
    async def chat_vision(
        self,
        message: str,
        image: Union[str, Path],
        **kwargs
    ) -> str
    
    def get_memory_stats(self) -> Dict[str, float]
    
    async def cleanup(self) -> None
\end{lstlisting}

\subsubsection{InferenceEngine}

The unified inference backend abstraction:

\begin{lstlisting}[language=Python, caption=InferenceEngine Interface]
class InferenceEngine:
    def __init__(
        self,
        model_path: Union[str, Path],
        backend: str = "llama_cpp",
        config: Optional[Dict[str, Any]] = None,
    )
    
    async def initialize(self) -> None
    
    async def generate(
        self,
        prompt: str,
        max_tokens: int = 512,
        temperature: float = 1.0,
        top_p: float = 0.95,
        top_k: int = 40,
        stream: bool = False,
    ) -> Union[str, AsyncIterator[str]]
    
    async def chat(
        self,
        messages: List[Dict[str, str]],
        **generation_params
    ) -> Union[str, AsyncIterator[str]]
    
    async def cleanup(self) -> None
\end{lstlisting}

\subsubsection{MemoryManager}

The memory management subsystem tracks allocations and enforces limits:

\begin{lstlisting}[language=Python, caption=MemoryManager Key Methods]
class MemoryManager:
    # Hardware limits
    TOTAL_RAM_GB = 28.0
    RESERVE_OS_GB = 4.0
    MAX_AVAILABLE_GB = 24.0
    SAFETY_FACTOR = 0.85
    
    def select_quantization(
        self,
        model_params_billions: float
    ) -> str
    
    def calculate_max_context(
        self,
        model_quantized_gb: float,
        kv_cache_type: str = "turbo3",
    ) -> int
    
    def get_optimal_settings(
        self,
        model_params_billions: float
    ) -> Dict[str, Any]
    
    def check_memory_available(
        self,
        required_gb: float
    ) -> bool
\end{lstlisting}

\subsection{Async Architecture}

Miniforge's async-first design provides several advantages:

\begin{enumerate}
    \item \textbf{Non-blocking:} Multiple concurrent requests can be processed without blocking the event loop
    \item \textbf{Streaming:} Natural fit for streaming token generation
    \item \textbf{Resource Sharing:} Concurrent tasks can share model weights while having separate KV caches
    \item \textbf{Integration:} Compatible with modern async web frameworks (FastAPI, Sanic, etc.)
\end{enumerate}

The implementation uses Python's \texttt{asyncio} with thread pool delegation for CPU-intensive operations:

\begin{lstlisting}[language=Python, caption=Async Pattern in Miniforge]
async def generate(self, prompt: str, **kwargs) -> str:
    if not self._initialized:
        await self.initialize()
    
    # Run blocking backend call in thread pool
    loop = asyncio.get_event_loop()
    result = await loop.run_in_executor(
        None,  # Default executor
        lambda: self._backend.generate_sync(prompt, **kwargs)
    )
    return result
\end{lstlisting}

\subsection{Error Handling and Recovery}

The architecture includes robust error handling at multiple levels:

\begin{itemize}
    \item \textbf{Backend Fallback:} If llama.cpp fails to load, automatically falls back to transformers
    \item \textbf{Context Fallback:} If requested context exceeds available memory, gradually reduces context size until initialization succeeds
    \item \textbf{OOM Prevention:} MemoryManager rejects requests that would exceed available memory
    \item \textbf{Graceful Degradation:} Reduces batch size or quantization level under memory pressure
\end{itemize}

\subsection{Extension Points}

The architecture supports future extensions through well-defined interfaces:

\begin{itemize}
    \item \textbf{New Backends:} Implement the backend interface to add support for additional inference engines
    \item \textbf{Custom Tools:} Tool handlers can be arbitrary Python functions, sync or async
    \item \textbf{Vision Models:} VisionProcessor can be extended for different image encoding schemes
    \item \textbf{Streaming Handlers:} Custom streaming strategies can be implemented for different use cases
\end{itemize}

\subsection{Thread Safety and Concurrency}

Model weights are read-only after initialization, enabling safe concurrent access. However, the KV cache and generation state are per-request. The implementation uses:\

\begin{itemize}
    \item \texttt{asyncio.Lock} for exclusive access during generation
    \item Thread-local storage for request-specific state
    \item Immutable model configuration shared across requests
\end{itemize}

This design enables multiple concurrent generations while maintaining correctness.
